\documentclass[11pt,a4paper]{article}
\usepackage{xeCJK}
\usepackage{fontspec}
\usepackage{xcolor}
\usepackage{fancyhdr}
\usepackage{booktabs}
\usepackage{array}
\usepackage{amssymb}
\usepackage{geometry}
\usepackage[protrusion=true]{microtype}
\geometry{margin=1.8cm,top=2cm,headheight=15pt}
\linespread{1.05}

\setmainfont{Baskerville}
\setCJKmainfont{Songti SC}
\newfontfamily{\starfont}{Songti SC}

\definecolor{wrongred}{RGB}{198,40,40}
\definecolor{wrongbg}{RGB}{253,235,233}
\definecolor{rulegray}{RGB}{200,200,200}

\setlength{\emergencystretch}{3em}
\setlength{\fboxsep}{3pt}
\setlength{\parindent}{0pt}

\newcommand{\blank}{\underline{\hspace{2.8cm}}}
\newcommand{\checkmarkbox}{\fbox{\rule{0pt}{0.8ex}\rule{0.8ex}{0pt}}}
\newcommand{\STAR}{\textcolor{wrongred}{\starfont ★}}
% 原卷作答（红色）
\newcommand{\W}[1]{\textcolor{wrongred}{#1}}
% 标准答案（加粗）
\newcommand{\A}[1]{\textbf{#1}}
% 判定标记
\newcommand{\OK}{\textcolor{wrongred}{\ensuremath{\checkmark}}}
\newcommand{\NO}{\textcolor{wrongred}{$\times$}}
% 表格正文列：左对齐，避免 p 列两端对齐造成的 Underfull
\newcolumntype{L}[1]{>{\raggedright\arraybackslash}p{#1}}

\pagestyle{fancy}
\fancyhf{}
\fancyhead[L]{\footnotesize 英语短语默写 小蓝 P21--30（错题订正）}
\fancyhead[R]{\footnotesize 赵文瀚}
\fancyfoot[C]{\thepage}
\renewcommand{\headrulewidth}{0.4pt}
\renewcommand{\footrulewidth}{0pt}

\begin{document}

\begin{center}
{\LARGE\bfseries 英语短语默写 小蓝 P21--30}\\[0.25em]
{\large 错题订正与复盘}\\[0.35em]
{\footnotesize 2029 届高一英语 · 小蓝短语默写 P21--30（共 25 条）· 2026-10}
\end{center}

\vspace{-0.2em}
\noindent{\color{rulegray}\rule{\textwidth}{0.8pt}}

\vspace{0.5em}
\noindent\fcolorbox{black!15}{black!4}{%
  \parbox{\dimexpr\linewidth-2\fboxsep-2\fboxrule\relax}{%
    \small 日期：2026-10\qquad 来源：小蓝短语默写 P21--30（共 25 条）\qquad 姓名：赵文瀚\\[0.3em]
    判错：\textcolor{wrongred}{\textbf{16}}\,/\enspace 25\qquad
    判对：\textbf{9}\,/\enspace 25\qquad
    重做日期：\underline{\hspace{2.4cm}}}}

\vspace{0.6em}
\noindent{\footnotesize 说明：下表按原卷题号列全 25 条。\W{红色}＝原卷作答，\A{加粗}＝标准答案；\OK\ 表示判对，\NO\ 表示判错。\STAR\ 表示该题暴露出成系统的搭配/词形问题，需优先攻破。\textbf{注意：本卷照片中无可识别的红笔批改痕迹，判定与标准答案由 AI 按标准英语给出，请对照原卷核对后再二刷。}}

\vspace{0.4em}
\renewcommand{\arraystretch}{0.96}
\noindent{\footnotesize\begin{tabular}{@{}L{0.72cm}@{\hspace{0.08cm}}L{2.03cm}@{\hspace{0.08cm}}L{1.45cm}@{\hspace{0.08cm}}L{4.6cm}@{\hspace{0.08cm}}L{4.7cm}@{\hspace{0.08cm}}L{0.75cm}@{\hspace{0.08cm}}L{1.55cm}@{\hspace{0.08cm}}L{0.9cm}@{}}
\toprule
\textbf{题号} & \textbf{中文短语} & \textbf{提示词} & \textbf{我的作答} & \textbf{标准答案} & \textbf{判定} & \textbf{错误类型} & \textbf{二刷} \\
\midrule
1  & 除了做某事别无选择     & alternative & have no alternative but to do & — & \OK & — & \checkmarkbox \\
2  & 业余自行车手           & —           & \W{amateur bicycler} & \A{amateur cyclist} & \NO & 用词错误 & \checkmarkbox \\
3  & …及其他东西            & among       & among other things & — & \OK & — & \checkmarkbox \\
4  & 拿游戏取乐             & amuse       & \W{amuse with games} & \A{amuse oneself with games} & \NO & 漏词 & \checkmarkbox \\
5  & 数据分析               & —           & analyse data & — & \OK & — & \checkmarkbox \\
6  & 对拖延症感到气愤 \STAR & —           & \W{be angry \textbf{with} the \textbf{delay}} & \A{be angry \textbf{about procrastination}} & \NO & 用词错误 & \checkmarkbox \\
7  & 对…负责 \STAR          & answer      & \W{take answer to} & \A{answer for} & \NO & 固定搭配 & \checkmarkbox \\
8  & 一个又一个胜利         & —           & \W{one \textbf{winning} after another} & \A{one \textbf{victory} after another} & \NO & 用词错误 & \checkmarkbox \\
9  & 带某人参观             & around      & show sb. around & — & \OK & — & \checkmarkbox \\
10 & 预料明年的销售量增加   & —           & anticipate the rise in sales next year & — & \OK & — & \checkmarkbox \\
11 & 对…渴望 \STAR          & anxious     & \W{be anxious \textbf{about} …（另写 can't wait for）} & \A{be anxious \textbf{for} …（be desperate for）} & \NO & 介词搭配 & \checkmarkbox \\
12 & 掌声四起 \STAR         & —           & \W{applaud／applause} & \A{burst into applause} & \NO & 固定搭配 & \checkmarkbox \\
13 & 在我看来 \STAR         & appear      & \W{in my appearance} & \A{it appears to me that} & \NO & 用词错误 & \checkmarkbox \\
14 & 要求道歉 \STAR         & —           & \W{be wanted to apologize} & \A{require an apology} & \NO & 固定搭配 & \checkmarkbox \\
15 & 为某事向某人道歉       & apologize   & apologize to sb. for sth. & — & \OK & — & \checkmarkbox \\
16 & 政府呼吁人人节约用水 \STAR & appeal  & \W{The government appeals everyone to save water} & \A{appeals \textbf{to} everyone to save water} & \NO & 漏词 & \checkmarkbox \\
17 & 你有兴趣去合资公司上班吗？ \STAR & appeal  & \W{Do you have interest in working in an …} & \A{Does the idea of working for a joint venture company appeal to you?} & \NO & 固定搭配 & \checkmarkbox \\
18 & 下结论 \STAR           & arrive      & \W{arrive\textbf{d} a conclusion} & \A{arrive \textbf{at} a conclusion} & \NO & 漏词 & \checkmarkbox \\
19 & 请假                   & apply       & \W{apply for a leaving} & \A{apply for leave} & \NO & 词形＋多词 & \checkmarkbox \\
20 & 派某人任某职           & —           & appoint sb. to a post & — & \OK & — & \checkmarkbox \\
21 & 如果…我将不胜感激 \STAR & appreciate & \W{I'll appreciate it if that} & \A{I'd appreciate it very much if …} & \NO & 固定搭配 & \checkmarkbox \\
22 & 对…欣赏               & apprecia\-tion & show appreciation for … & — & \OK & — & \checkmarkbox \\
23 & 不同意我辍学 \STAR     & approve     & \W{don't approve \textbf{for} my getting out of school} & \A{don't approve \textbf{of} my \textbf{dropping out of} school} & \NO & 用词错误 & \checkmarkbox \\
24 & 拱腰                   & —           & arch one's back & — & \OK & — & \checkmarkbox \\
25 & 安排医生给他看病 \STAR & arrange     & \W{arrange a doctor to took after him} & \A{arrange \textbf{for} a doctor to \textbf{see} him} & \NO & 固定搭配 & \checkmarkbox \\
\bottomrule
\end{tabular}}

\vspace{0.3em}
\noindent{\footnotesize\textbf{原卷旁注}：（一）本卷照片偏暗、色彩饱和度极低，未见可识别的红笔痕迹，转录以 OCR 识别结果为准，涂改与划线无法确证；请对照 {\xeCJKsetup{CJKecglue={}}\texttt{phrase\_02\_p21\_30\_01\_原卷\_1.jpg}} 复核后再二刷。（二）第 6 题作答行上方另有 \emph{procrastination} 字样，第 11 题上方有 \emph{be desperate for}，第 23 题上方有 \emph{dropping}，第 18 题作答读作 \emph{arrived a conclusion}（\emph{arrive} 后有一字母，多读作 d）；本次全页 OCR 未检出任何 \emph{at}／\emph{for} 行间补词，故「下方补 at / for」不作定论，但两题作答均缺该介词，判错不因此改变。（三）第 7、14 题笔迹潦草，OCR 读作 \emph{take answer to} 与 \emph{be wanted to apologize}，按标准答 \emph{answer for}、\emph{require an apology} 处理。（四）第 17 题学生写 \emph{Do you have interest in working in an …}，未用提示词 appeal 的考点句型，故判错；标准答 \emph{Does the idea of working for a joint venture company appeal to you?}（\emph{sth. appeal to sb.}）。原卷未见可识别划痕，若与老师批改不符请再改。（五）第 22 题 \emph{to/for}、第 24 题记号存疑，暂判对。（六）2026-10 已用《2026 年上海市高考英语词汇用法手册》复核：9 条在该手册「词组」栏检到且一致，其余未收；无冲突。}

\newpage
\begin{center}
{\Large\bfseries 错因归类与复盘}
\end{center}
\vspace{-0.4em}
\noindent{\color{rulegray}\rule{\textwidth}{0.8pt}}

\vspace{0.5em}
\noindent 说明：按错误性质归类 16 条错题，供汇总对账；逐条二刷结果见第 1 页「二刷」列。

\vspace{0.5em}
\renewcommand{\arraystretch}{1.35}
\begin{center}
\begin{tabular}{@{}p{4.6cm} p{10.3cm} c@{}}
\toprule
\textbf{错误类型} & \textbf{题号} & \textbf{条数} \\
\midrule
固定搭配错误 & {\small 7, 12, 14, 17, 21, 25} & 6 \\
用词错误（近义误用） & {\small 2, 6, 8, 13, 23} & 5 \\
漏词 & {\small 4, 16, 18} & 3 \\
词形＋多词（leaving$\to$leave、多 a） & {\small 19} & 1 \\
介词搭配（about / for） & {\small 11} & 1 \\
\midrule
\textbf{合计} & \textbf{16 条} & \textbf{16} \\
\bottomrule
\end{tabular}
\end{center}

\vspace{0.5em}
\noindent\textbf{关联知识点标签}
\vspace{0.25em}

\noindent\begin{tabular}{@{}p{8.2cm}@{\hspace{0.4cm}}p{8.2cm}@{}}
\begin{minipage}[t]{\linewidth}
\small
\STAR\ \textbf{appeal 两用}：appeal \textbf{to} sb. to do sth.（呼吁）；sth. appeal \textbf{to} sb.（吸引）\\[0.2em]
\STAR\ \textbf{approve}：approve \textbf{of}（doing）sth.；\textbf{drop out of} school（辍学，不是 get out of school）\\[0.2em]
\STAR\ \textbf{appreciate 句式}：I \textbf{would} appreciate it very much \textbf{if} …（不用 I'll、不加 that）\\[0.2em]
\STAR\ \textbf{arrange}：arrange \textbf{for} sb. \textbf{to do} sth.（for 不可省）
\end{minipage} &
\begin{minipage}[t]{\linewidth}
\small
\STAR\ \textbf{固定搭配}：answer \textbf{for}（负责）、burst into applause（掌声四起）、require an apology（道歉）、arrive \textbf{at} a conclusion\\[0.2em]
\STAR\ \textbf{介词搭配}：be anxious \textbf{for}（渴望），不是 about；be angry \textbf{about} sth.\\[0.2em]
\STAR\ \textbf{词形/用词}：\textbf{cyclist}（非 bicycler）、\textbf{leave}（名词，非 leaving）、\textbf{victory}（非 winning）\\[0.2em]
\STAR\ \textbf{范围}：小蓝短语 P21--30（a- 家族：alternative / amuse / answer / anxious / appear / apologize / appeal / arrive / apply / appreciate / approve / arrange；全卷 25 条按字母序，\STAR\ 题为待攻破项）
\end{minipage} \\
\end{tabular}

\vspace{0.65em}
\noindent\fcolorbox{black!15}{black!4}{%
  \parbox{\dimexpr\linewidth-2\fboxsep-2\fboxrule\relax}{%
    \textbf{复盘记录}\\[0.35em]
    最薄弱的板块：\underline{\hspace{9.0cm}}\\[0.55em]
    主要错误原因：\checkmarkbox 固定搭配不牢\quad \checkmarkbox 介词记混\quad \checkmarkbox 漏词/多词\quad \checkmarkbox 用词/拼写\quad \checkmarkbox 空白未答\\[0.55em]
    本次复习的改进措施：\underline{\hspace{10.1cm}}\\[0.55em]
    \hspace{3.2em}\underline{\hspace{10.1cm}}\\[0.55em]
    下次复习日期：\underline{\hspace{2.2cm}}\qquad
    当前掌握度：\checkmarkbox 仍薄弱\quad \checkmarkbox 基本掌握\quad \checkmarkbox 已掌握
  }}

\end{document}
